\chapter{Memory Objects --- Ownership That Crosses Boundaries}

A memory object answers: \textbf{``Where is the data, and who currently owns
it?''} It is the mechanism by which data crosses protection-domain boundaries
without copying bytes. This chapter explains CharlotteOS's memory-object model.

\section{The problem with copying}

In a conventional OS, moving data between processes typically means copying:
the kernel reads from the sender's address space and writes to the receiver's.
For large payloads (packet buffers, file blocks, IPC messages), this copy
dominates the cost of the operation.

The alternative --- shared memory --- eliminates the copy but eliminates the
safety. Two processes with writable mappings to the same physical pages can
corrupt each other's data. Synchronization is advisory; enforcement is absent.

CharlotteOS adopts a third approach: \textbf{ownership transfer through page-table
manipulation.} When a memory object moves from one domain to another, the kernel
changes the page-table mappings. The sender loses access; the receiver gains it.
No bytes are copied. The MMU enforces the ownership.

\section{First-class memory objects}

A memory object is a dedicated kernel object representing page-backed storage with
explicit ownership and capability semantics. The kernel manages:

\begin{itemize}
    \item \textbf{Physical frames} --- the backing storage.
    \item \textbf{Mappings} --- which domains have virtual addresses pointing to
        the pages, and with what permissions.
    \item \textbf{Ownership} --- which domain currently holds the master
        capability.
    \item \textbf{Transfer state} --- whether the object is idle, in transit,
        borrowed, or undergoing DMA.
    \item \textbf{Lifetime} --- when the object is freed, all mappings are
        revoked and all capabilities to it are invalidated.
\end{itemize}

Userspace never exchanges raw pointers across protection domains. Instead, IPC
messages attach \textbf{memory-object capabilities} with declared transfer modes.

\section{Four transfer modes}

CharlotteOS defines four transfer modes, enforced by the MMU:

\subsection{Copy}

The receiver gets a byte-for-byte duplicate. The sender retains the original,
unmodified.

\begin{itemize}
    \item \textbf{Use when:} the data is small, the sender needs to keep using it,
        or the receiver is not trusted with the original.
    \item \textbf{Kernel work:} allocate new physical frames, copy the contents,
        map them into the receiver's address space.
    \item \textbf{Semantics:} simplest; both domains can read their respective
        copies independently.
\end{itemize}

\subsection{Move}

Ownership transfers to the receiver. The sender loses all mappings and all
capability rights. The receiver becomes the new owner, responsible for subsequent
transfer or destruction.

\begin{itemize}
    \item \textbf{Use when:} the sender is finished with the data and wants to
        hand it off without copying.
    \item \textbf{Kernel work:} validate sender ownership, prepare receiver
        mappings, remove sender mappings, perform TLB invalidation, commit
        ownership, publish the IPC message.
    \item \textbf{Semantics:} physical pages never move. Mappings change. After
        the move, any attempt by the sender to access the (now unmapped) virtual
        addresses faults.
\end{itemize}

\subsection{BorrowRead}

The receiver gets a temporary read-only mapping. The sender retains ownership.
The borrow is revoked when the receiver replies (or the call is cancelled, or
the server dies).

\begin{itemize}
    \item \textbf{Use when:} the receiver needs to inspect data but must not
        modify it.
    \item \textbf{Kernel work:} map the pages into the receiver's address space
        with read-only permissions. Track the borrow so it can be revoked.
    \item \textbf{Semantics:} writable aliases in the receiver are forbidden by
        the page-table permissions. The receiver cannot write even through unsafe
        code. Revocation at reply time unmap the pages.
\end{itemize}

\subsection{BorrowWrite}

The receiver gets a temporary exclusive writable mapping. The sender's access is
revoked for the duration of the borrow. When the receiver replies, sender access
is restored.

\begin{itemize}
    \item \textbf{Use when:} the receiver needs to write into a buffer the sender
        will later read (e.g., a read-into-user-buffer pattern). This is
        \emph{mutable lending}.
    \item \textbf{Kernel work:} map the pages writable into the receiver,
        temporarily revoke (unmap or read-protect) sender access, track the
        borrow for reply-time revocation.
    \item \textbf{Semantics:} the hardware enforces mutual exclusion. The sender
        cannot access the pages while they are lended out. No amount of unsafe
        Rust can violate this --- the MMU will fault.
\end{itemize}

\section{Why hardware enforcement matters}

The transfer modes are enforced by page tables, not by Rust's type system:

\begin{itemize}
    \item A C program running in a protection domain cannot violate the
        constraints any more than a Rust program can.
    \item A buffer overrun cannot read pages that are not mapped.
    \item An unsafe block cannot write to a \texttt{BorrowRead} buffer.
    \item A use-after-drop cannot resurrect a moved memory object.
\end{itemize}

Rust's ownership complements the kernel enforcement --- it catches errors within a
domain at compile time. But the kernel is the ultimate guarantor for cross-domain
safety, so non-Rust or buggy-Rust processes are still contained by page tables and
capability checks.

\section{Control plane versus data plane}

Endpoint IPC belongs to the \textbf{control plane}: it configures connections,
carries opcodes, sets up device queues.

Memory objects belong to the \textbf{data plane}: they carry the actual bytes
(packets, blocks, buffers, strings).

For networking specifically:

\begin{itemize}
    \item \textbf{Control plane:} endpoint IPC configures NIC queues, sets MTU,
        registers buffer pools.
    \item \textbf{Data plane:} descriptor rings coordinate producer/consumer
        state; packet payloads move as memory objects.
\end{itemize}

This separation avoids turning every packet into a copied IPC payload while
preserving strict ownership.

\section{DMA as a third ownership dimension}

When a device performs DMA, a memory object enters a third ownership state:

\begin{itemize}
    \item \textbf{CPU ownership:} the CPU reads and writes the pages normally.
    \item \textbf{DMA ownership:} the device is reading or writing the pages; the
        CPU must not access them.
    \item \textbf{Rust/API ownership:} which domain holds the capability.
\end{itemize}

The architecture calls for tracking DMA state per memory object and removing
CPU access while a device owns a buffer. The current prototype exposes physical
addresses to trusted userspace drivers and does not yet implement IOMMU/SMMU
isolation. Consequently, capability isolation does not currently constrain a
misbehaving DMA-capable device or driver.

\section{Why memory objects matter for critical software}

\begin{enumerate}
    \item \textbf{Page transfer without payload copying.} For aligned,
        page-backed objects, ownership can move by changing mappings rather than
        copying bytes. Whether this is faster depends on payload size and
        mapping/TLB costs.

    \item \textbf{Cross-domain stale access is constrained by mappings.}
        When a memory object is moved, the sender's mappings are removed. Any
        later access through those mappings faults. When a borrow is revoked,
        the borrower's mappings are removed. This does not rule out bugs in
        unsafe code, page-table maintenance, or DMA.

    \item \textbf{Borrows have defined lifetimes.} A borrow is always tied to an
        IPC call. When the call completes (reply, cancellation, timeout, or
        server death), the borrow ends. There is no ``I forgot to return it.''

    \item \textbf{DMA ownership is an architectural requirement.} Full
        kernel-tracked pinning and IOMMU enforcement remain future work; current
        drivers are trusted to retain buffers until hardware completion.

    \item \textbf{Hardware enforcement, not language enforcement.} The page tables
        are the ultimate arbiter. No amount of unsafe code, C FFI, or compiler
        bug should access CPU pages that are not mapped, assuming correct MMU
        setup and TLB maintenance. DMA requires the not-yet-implemented IOMMU
        path for an equivalent boundary.
\end{enumerate}

\endinput
